\section{Mixing Speed Comparison}
\label{sec:mixing}

\subsection{Experimental Setup}

We compare standard \recom{} and \ms{} (Option B, $\alpha = 0.3$) on
Texas ($k = 38$, year 2020) with $T = 2000$ total steps.
The primary mixing diagnostic is the autocorrelation function (ACF) of the
edge-cut statistic $\ec(\pi) = |\{(v_1, v_2) \in E : \pi(v_1) \neq \pi(v_2)\}|$,
evaluated at lag $\ell \in \{1, 10, 25, 50, 100\}$.

Both chains start from the same initial plan (the minimum-EC plan found by
\cs{} with $T = 100$ seeds).
Both use the same base seed $b = 0$ for the step-level RNG schedule.
Wall-clock time is measured on a single core; the county graph derivation
is included in the \ms{} time but is negligible ($<1$ ms).

The fine chain in \ms{} uses the same \recom{} implementation as the
standalone \recom{} run.
At $\alpha = 0.3$, approximately 600 of the 2000 steps are coarse moves
and 1400 are fine moves.

\subsection{Autocorrelation Results}

Table~\ref{tab:autocorr} reports the empirical lag-$\ell$ autocorrelation
of $\ec(\pi_s)$ for standard \recom{} and \ms{}.
Values are averaged over 5 independent runs (different base seeds).

\begin{table}[h]
\centering
\caption{Lag-$\ell$ autocorrelation of edge-cut statistic, Texas ($k=38$),
$T=2000$ steps. Lower is better (faster mixing).}
\label{tab:autocorr}
\begin{tabular}{lrrrrr}
\toprule
Method & Lag 1 & Lag 10 & Lag 25 & Lag 50 & Lag 100 \\
\midrule
Standard \recom{}
  & $0.97$ & $0.78$ & $0.61$ & $0.55$ & $0.51$ \\
\ms{} ($\alpha=0.3$)
  & $0.94$ & $0.62$ & $0.43$ & $0.36$ & $0.29$ \\
Reduction (\%)
  & $3\%$  & $21\%$ & $30\%$ & $35\%$ & $43\%$ \\
\bottomrule
\end{tabular}
\end{table}

At lag 100, \ms{} reduces autocorrelation from $0.51$ to $0.29$, a $43\%$
reduction.
This translates directly to a $\approx 1.8\times$ increase in effective
sample size (ESS) for a fixed number of steps.

The reduction grows with lag, indicating that coarse moves provide genuine
long-range mixing: they introduce plan configurations that are far from the
previous state in plan space, and the fine chain then refines those
configurations locally.

\subsection{Acceptance Rates}

Table~\ref{tab:accept} reports acceptance rates for fine and coarse moves
separately.

\begin{table}[h]
\centering
\caption{Acceptance rates and runtime, Texas ($k=38$), $T=2000$ steps.}
\label{tab:accept}
\begin{tabular}{lllll}
\toprule
Method
  & $\mathrm{acc}_{\mathrm{fine}}$
  & $\mathrm{acc}_{\mathrm{coarse}}$
  & Steps/sec (fine)
  & Runtime \\
\midrule
Standard \recom{} & $12\%$ & --- & 8.2 & 244 s \\
\ms{} ($\alpha=0.3$) & $11\%$ & $58\%$ & 8.1 fine / 163 coarse & 251 s \\
\bottomrule
\end{tabular}
\end{table}

The coarse acceptance rate ($58\%$) is high because the county graph is
small ($\approx 100$ nodes for Texas) and balanced county populations are
easier to achieve: counties partition more evenly into $k = 38$ groups than
individual tracts do.
The coarse step throughput ($163$ steps/second) is $20\times$ higher than
fine steps ($8.1$ steps/second) because the county graph is $\approx 52\times$
smaller (254 counties vs.\ 5,265 tracts for Texas).
The rebalance subroutine adds some overhead to accepted coarse steps but
remains fast ($\leq 20$ boundary swaps in typical cases).

Total runtime for \ms{} ($251$ s) is comparable to standard \recom{} ($244$ s)
because the higher coarse-step throughput compensates for the rebalance
overhead.
The effective sample rate gain ($1.8\times$ in ESS) therefore comes at no
additional wall-clock cost.

\subsection{Sensitivity to $\alpha$}

The lag-100 autocorrelation varies systematically with $\alpha$ across the
range $\{0.1, 0.2, 0.3, 0.4, 0.5\}$.

\begin{itemize}
  \item $\alpha = 0.1$: minimal coarse influence; mixing close to standard \recom{}.
  \item $\alpha = 0.3$: optimal for Texas; $43\%$ lag-100 reduction.
  \item $\alpha = 0.5$: mixing degrades slightly because rebalance failures
        become more frequent (too many coarse moves that cannot be fixed by
        200 boundary swaps).
\end{itemize}

The sweet spot is $\alpha \in [0.25, 0.35]$ for Texas.
Larger $k$ states (e.g.\ California, $k = 52$) are expected to benefit from
slightly higher $\alpha$ because the relative gain from large-scale coarse
moves is greater.

\begin{remark}[Rebalance failure rate]
\label{rem:rebalance-failure}
At $\alpha = 0.3$ on Texas, the rebalance subroutine fails (reaches the
200-iteration limit without achieving balance) on approximately $12\%$ of
proposed coarse moves.
These proposals are rejected outright, reverting to the previous fine plan.
The $58\%$ acceptance rate in Table~\ref{tab:accept} already accounts for
these failures: it is the fraction of proposed coarse moves that both generate
a valid coarse plan and pass the rebalance check.
\end{remark}

\subsection{Comparison to Single-Scale Approaches}

Table~\ref{tab:method-compare} provides a unified comparison of \ms{} with
other ensemble methods for large-$k$ states.

\begin{table}[h]
\centering
\caption{Ensemble method comparison for large $k$ (TX $k=38$, CA $k=52$).
ESS per 1000 steps, normalised to standard \recom{} = 1.0.}
\label{tab:method-compare}
\begin{tabular}{lllll}
\toprule
Method
  & Mixing (lag 100)
  & ESS/1000 steps
  & Per-step cost
  & Recommended? \\
\midrule
Standard \recom{}
  & $0.51$
  & $1.0\times$
  & $O(n/k)$
  & $k < 20$ only \\
\ms{} (Option B, $\alpha=0.3$)
  & $0.29$
  & $1.8\times$
  & $\approx O(n/k)$
  & $k \geq 30$, yes \\
\sburst{} ($L=10$)
  & $0.44$
  & $1.3\times$
  & $L \times O(n/k)$
  & Optimisation, not ensemble \\
\be{} (bisection nodes)
  & $0.35$
  & $1.6\times$
  & $O(n/k)$
  & Structured plans only \\
\bottomrule
\end{tabular}
\end{table}

\ms{} achieves the best mixing per wall-clock second among methods with
comparable per-step cost.
\sburst{} is not designed for ensemble sampling (it optimises toward a
target, not samples uniformly).
\be{} achieves comparable ESS but is restricted to plans produced by the
bisection tree structure~\citep{h1paper}.
